\documentclass{naklatex}
\usepackage[utf8]{inputenc}
\usepackage[T1]{fontenc}
\usepackage{geometry}
\usepackage{graphicx}
\usepackage{float}
\usepackage{hyperref}
\usepackage[english,indonesian]{babel}

\title{Optimasi Pemrosesan Template LaTeX Menggunakan Pendekatan Berbasis AI}
\author{Academic Support Agent}
\date{\today}

\begin{document}
\maketitle

\begin{abstract}

\end{abstract}

\chapter{PENDAHULUAN}


\chapter{TINJAUAN PUSTAKA}


% Original content (commented out):
% \documentclass[]{naklatex} 
\usepackage{lipsum}

% Zitationseinstellungen
\usepackage[backend=biber, style=apa]{biblatex}
\ExecuteBibliographyOptions{sorting=nty}
\addbibresource{quellen.bib}


\begin{document}

    %% In der NAKLATEX Klasse sind unten noch weitere händisch erstellte Befehle
    %% wie "\deckblatt". Dort kann nachgeschaut werden, was bereits implementiert
    %% wurde und Anpassungen vornehmen.

    %% Deckblatt, Inhaltsverzeichnis und Co
    \begin{frontmatter}
        \deckblatt
        \inhaltsverzeichnis
    \end{frontmatter}
    
    %% Kapitel
    \begin{mainmatter}
        \input{chapter/Einleitung}
    \end{mainmatter}

    \begin{backmatter}
        \literaturverzeichnis
    \end{backmatter}

\end{document}

\end{document}